\chapter*{Executive Summary}
\addcontentsline{toc}{chapter}{Executive Summary}

This report documents the industrial attachment undertaken at \textbf{Texlab IT}, Rajshahi, Bangladesh, from February 24, 2026 to March 10, 2026. Texlab IT is a prominent digital agency and IT training institute specializing in web services, full-stack application development, and skill training. During the two-week attachment, the primary objective was to design, develop, and stabilize a production-grade web application called the \textbf{URL Shrinker API} under industrial supervision, bridging the gap between academic theory and practical software engineering workflows.

The URL Shrinker API is a high-performance URL shortening service designed to handle substantial request throughput. It comprises a Go-based backend and a React/Next.js-based frontend. The system includes JWT-based user authentication, custom short code generation, public redirects with high-speed caching using a Redis-based cache-aside strategy, request rate limiting, and an analytics engine that tracks redirect statistics. It also features a background worker routine that runs periodically to remove expired links from the database. 

During the internship, several engineering challenges were addressed and resolved, including Turbopack compilation instability on the frontend, database pool synchronization utilizing PostgreSQL and the `sqlc` library, cache stampede prevention, and handling CORS and authentication middleware in Go. The development process adhered to professional software engineering standards, utilizing Docker for local containerization and GitHub for version control.

This attachment facilitated significant professional and technical growth, enhancing skills in Go database programming, cache-aside system design, frontend state management, and debugging in a fast-paced development environment. Furthermore, it provided exposure to industrial practices in safety, legal compliance (GDPR and data privacy), professional ethics, and collaborative development workflows, laying a solid foundation for lifelong learning in engineering.
